\documentclass[11pt]{article}
\usepackage{setspace}
\usepackage[a4paper,margin=1in]{geometry}
\geometry{margin=0.85in}
\setstretch{1.0}
\usepackage{graphicx}
\usepackage{array}
\usepackage{setspace}
\usepackage{parskip}
\usepackage{ragged2e}

\pagestyle{empty}
\setstretch{1.05}

\begin{document}


% Title
\begin{center}
    {\Large \textbf{Job Application}} \\[3cm]
    
    {\Huge \textbf{Zhongli Wang}} \\[0.4cm]
    
    {\LARGE Application for ET Consultant - Geospatial Data Scientist}
\end{center}

\vspace{0.6cm}

% Large picture in the middle
\begin{center}
    \IfFileExists{../zhongli-wang-profile.jpg}{
        \includegraphics[
            width=0.55\textwidth,
            height=0.45\textheight,
            keepaspectratio
        ]{../zhongli-wang-profile.jpg}
    }{
        \vspace{0.45\textheight}
    }
\end{center}

\vspace{0.7cm}

% \vspace{3cm}


% Personal details
\begin{center}
\renewcommand{\arraystretch}{1.35}
\begin{tabular}{>{\bfseries}p{4.2cm} p{8.5cm}}


Address: & Astastrasse 22, 33617 Bielefeld, Germany \\

Number: & +49 160 9541 6015 \\

Email: & nkwangzhongli@gmail.com \\
\end{tabular}
\end{center}

\vfill

% \begin{center}
%     \small Last updated: \today
% \end{center}
\newpage
\enlargethispage{4\baselineskip}

\begin{flushright}
Zhongli Wang \\
Astastrasse 22 \\
33617 Bielefeld\\
Germany \\
\end{flushright}

\vspace{0.1cm}

\vspace{0.1cm}

\noindent
World Bank Group \\
Infrastructure Vertical - Disaster Management \\
Washington, DC, United States

\vspace{0.12cm}

\noindent
\textbf{Application for ET Consultant - Geospatial Data Scientist}

\vspace{0.1cm}

\noindent Dear World Bank Group recruiting team,

I am writing to express my strong interest in the ET Consultant - Geospatial Data Scientist position supporting the Global Program on Nature-Based Solutions for Climate Resilience. 
I am completing a Ph.D. in Economics at Bielefeld University, with my defense scheduled for September 21, 2026, 
and I hold an M.A. in Business and Economics with a major in Banking and Finance and a minor in Data Science. 
My background combines advanced Python programming, quantitative data analysis, computational modeling, and machine learning, 
providing a strong foundation for developing and applying geospatial analytics tools to support climate resilience and disaster risk management.

I am particularly drawn to the World Bank Group's work on Nature-Based Solutions and climate resilience in developing countries. 
The opportunity to develop and apply quantitative geospatial tools like the NBS Opportunity Scan, support data-driven decision-making for climate adaptation, 
and contribute to knowledge products that guide governments in implementing NBS aligns perfectly with my research approach: 
design computational frameworks, process and analyze large-scale datasets, validate model outputs, 
and translate complex analytical results into accessible insights for policymakers and practitioners.
My relevant experience comprises:

\begin{itemize}
    \item \textbf{Python-based data analysis and modeling:} In published research, I designed data collection and processing pipelines, implemented machine learning models (XGBoost, Random Forest, Neural Networks) using Python (NumPy, pandas, scikit-learn), and developed analytical frameworks for large-scale datasets—demonstrating transferable skills for geospatial data processing and environmental analytics.
    \item \textbf{Computational modeling and tool development:} I developed agent-based computational models in Python to analyze complex systems, spatial dynamics, and temporal patterns. This experience in building analytical tools, automating workflows, and validating model outputs is directly applicable to developing and refining geospatial analytics tools like the NBS Opportunity Scan.
    \item \textbf{Quantitative analysis and knowledge products:} My doctoral research involves parameter estimation, statistical analysis, and translating complex computational results into peer-reviewed publications and conference presentations. I have demonstrated ability to produce rigorous analytical work and communicate findings to academic and policy audiences.
    \item \textbf{International collaboration and communication:} I have presented research at international conferences across Europe and China, developed teaching materials, and worked in multicultural environments. I am fluent in English and German, with native Mandarin proficiency—valuable for engaging with diverse stakeholders and government counterparts in developing countries.
\end{itemize}

While my background is in computational economics rather than traditional geospatial science, my strong Python programming skills, experience with large-scale data analysis, and proven ability to quickly master new technical tools position me well to contribute to the NBS program. 
I am highly motivated to apply my quantitative and computational expertise to climate resilience and disaster risk management, and I am eager to deepen my knowledge of geospatial analytics, Earth observation data, and environmental monitoring. 
I am ready to support the development of analytical tools, contribute to knowledge products, and help translate technical innovations into practical applications for World Bank task teams and government clients in developing countries.

Thank you for considering my application! 
I would welcome the opportunity to discuss how my computational modeling expertise and analytical skills could contribute to the World Bank Group's mission to advance Nature-Based Solutions and climate resilience in developing countries.

\vspace{0.08cm}
\noindent Kind regards, \\
\noindent Zhongli Wang

% \vspace{0.1cm}
% \noindent\includegraphics[width=2.7cm,keepaspectratio]{../Wang-sig.jpg}

\end{document}
